% BHU PYQ -- Sociology: Introduction to Sociology-I (2024-25 NEP)
\documentclass[11pt,a4paper]{article}
\usepackage[a4paper,top=2.5cm,bottom=2.5cm,left=2.8cm,right=2.8cm,headheight=14pt]{geometry}
\usepackage{amsmath,amssymb}


\usepackage{fontspec}
\setmainfont{Kohinoor Devanagari}
\usepackage{microtype,fancyhdr,enumitem,hyperref,lastpage,parskip}

\hypersetup{colorlinks=true,urlcolor=black,linkcolor=black,pdftitle={SOCMJ11: Introduction to Sociology-I -- BHU NEP 2024-25}}

\pagestyle{fancy}\fancyhf{}
\renewcommand{\headrulewidth}{0.4pt}\renewcommand{\footrulewidth}{0.4pt}
\fancyhead[L]{\small \textit{Banaras Hindu University}}
\fancyhead[R]{\small \textit{SOCMJ11: Sociology-I (2024-25)}}
\fancyfoot[C]{\small Page \thepage\ of \pageref{LastPage}}

\newlist{parts}{enumerate}{2}
\setlist[parts,1]{label=(\alph*),leftmargin=*,itemsep=4pt,topsep=3pt}
\setlist[parts,2]{label=(\roman*),leftmargin=*,itemsep=2pt,topsep=2pt}

\newcommand{\pts}[1]{\hfill \textit{[#1]}}

\begin{document}

\begin{center}
  {\large\bfseries BANARAS HINDU UNIVERSITY}\\[2pt]
  {\normalsize B. A. (Hons.) Semester I Examination, 2024-25 (NEP)}\\[6pt]
  \rule{0.8\linewidth}{0.5pt}\\[6pt]
  {\large\bfseries SOCIOLOGY}\\[4pt]
  {\large\bfseries Paper Code: SOCMJ11 : Introduction to Sociology-I}\\[4pt]
  {\normalsize Time: 2:30 Hours \hfill Full Marks: 70}
\end{center}
\rule{\linewidth}{0.4pt}
\smallskip

\noindent \textit{Instructions: This question paper comprises three Sections A, B and C. All questions are compulsory. Write your Roll No.\ at the top immediately on receipt of this question paper.}

\smallskip
\noindent \textit{इस प्रश्न-पत्र में तीन खण्ड अ, ब तथा स हैं। सभी प्रश्न अनिवार्य हैं।}

\bigskip

\section*{SECTION -- A / खण्ड -- अ}
\noindent \textit{Note: Objective Type Questions. Each question carries 2 marks. Answer should not exceed 50 words.}\pts{5 $\times$ 2 = 10}

\noindent \textit{वस्तुनिष्ठ प्रश्न। प्रत्येक प्रश्न 2 अंकों का है। उत्तर 50 शब्दों से अधिक नहीं होना चाहिए।}

\begin{enumerate}[label=\textbf{\arabic*.},leftmargin=*,itemsep=8pt]
  \item 
  \begin{parts}
    \item Who is the father of Sociology?\pts{2}
    
    समाजशास्त्र के जनक कौन हैं? (Auguste Comte, 1838)

    \item Family is \underline{\hspace{3cm}}.\pts{2}
    
    परिवार है \underline{\hspace{3cm}}।
    \begin{parts}
      \item Association (समिति)
      \item Institution (संस्था)
      \item Both (दोनों)
    \end{parts}

    \item Who is the author of the book ``What is Sociology''?\pts{2}
    
    ``व्हाट इज सोशियोलॉजी'' पुस्तक के लेखक कौन हैं? (Alex Inkeles)

    \item Who said that, ``Society is a web of social relationships''?\pts{2}
    
    ``समाज सामाजिक संबंधों का जाल है'' यह किसने कहा है? (MacIver \& Page)

    \item Who gave the theory of `Class Struggle'?\pts{2}
    
    `वर्ग संघर्ष' का सिद्धान्त किसने दिया है? (Karl Marx)
  \end{parts}
\end{enumerate}

\bigskip

\section*{SECTION -- B / खण्ड -- ब}
\noindent \textit{Note: Short Answer Type Questions. Each question carries 10 marks. Answer all questions within 250 words.}\pts{3 $\times$ 10 = 30}

\noindent \textit{लघु उत्तरीय प्रश्न। प्रत्येक प्रश्न 10 अंकों का है। प्रत्येक प्रश्न का उत्तर 250 शब्दों से अधिक नहीं होना चाहिए।}

\begin{enumerate}[label=\textbf{\arabic*.},leftmargin=*,start=2,itemsep=12pt]

  \item Sociology is a scientific discipline. Explain it.\pts{10}
  
  समाजशास्त्र एक वैज्ञानिक विषय है। व्याख्या कीजिए।

  \begin{center}\textbf{OR / अथवा}\end{center}

  Define Community. Discuss its main characteristics.\pts{10}
  
  समुदाय को परिभाषित कीजिए। इसकी विशेषताओं की चर्चा कीजिए।

  \item Discuss the formalistic and synthetic schools of Sociology.\pts{10}
  
  समाजशास्त्र के स्वरूपात्मक तथा समन्वयात्मक सम्प्रदायों की चर्चा कीजिए।

  \begin{center}\textbf{OR / अथवा}\end{center}

  What is socialization? Discuss its importance in human life.\pts{10}
  
  समाजीकरण क्या है? इसके महत्त्व की चर्चा कीजिए।

  \item Briefly discuss the central ideas of C. Wright Mills' `Sociological Imagination'.\pts{10}
  
  सी० राइट मिल्स के `समाजशास्त्र की सम्भावना' के मुख्य विचारों की संक्षेप में चर्चा कीजिए।

  \begin{center}\textbf{OR / अथवा}\end{center}

  Define Social Group. Differentiate between primary and secondary groups.\pts{10}
  
  सामाजिक समूह को परिभाषित कीजिए। प्राथमिक एवं द्वितीयक समूह में अन्तर स्पष्ट कीजिए।

\end{enumerate}

\bigskip

\section*{SECTION -- C / खण्ड -- स}
\noindent \textit{Note: Long Answer Type Questions. Each question carries 15 marks. Answer all questions within 500 words.}\pts{2 $\times$ 15 = 30}

\noindent \textit{दीर्घ उत्तरीय प्रश्न। प्रत्येक प्रश्न 15 अंकों का है। उत्तर 500 शब्दों से अधिक नहीं होना चाहिए।}

\begin{enumerate}[label=\textbf{\arabic*.},leftmargin=*,start=5,itemsep=14pt]

  \item What are the major historical events that led to the emergence of Sociology as a discipline?\pts{15}
  
  वे कौन-सी प्रमुख ऐतिहासिक घटनाएँ थीं जिनके कारण समाजशास्त्र का एक विषय के रूप में उदय हुआ?

  \begin{center}\textbf{OR / अथवा}\end{center}

  Critically examine MacIver's statement: ``We belong to associations and not to institutions''.\pts{15}
  
  मेकाइवर के कथन: ``हम समितियों के सदस्य होते हैं, संस्थाओं के नहीं'' की समालोचनात्मक समीक्षा कीजिए।

  \item Discuss Auguste Comte's contributions to Sociology.\pts{15}
  
  समाजशास्त्र में अगस्ट कॉम्टे के योगदान की चर्चा कीजिए।

  \begin{center}\textbf{OR / अथवा}\end{center}

  What do you understand by Role and Status? Differentiate between Ascribed and Achieved Status.\pts{15}
  
  भूमिका एवं परिस्थिति से आप क्या समझते हैं? प्रदत्त एवं अर्जित परिस्थिति में अन्तर स्पष्ट कीजिए।

\end{enumerate}

\end{document}
